\chapter{Conclusion}

\section{Summary of Work}

% ─────────────────────────────────────────────────────────────────────────────

The efficient management of maritime traffic within heavily congested geographic bottlenecks, such as the Singapore Strait, necessitates an intelligent surveillance infrastructure. Relying solely on vessels that voluntarily broadcast their positional telemetry is insufficient for comprehensive situational awareness. This project successfully conceptualized and deployed a robust optical pipeline engineered to detect every active vessel within a camera's field of view. By persistently tracking unique ship identities through severe visual occlusions and algorithmically forecasting their future trajectories, the resulting architecture effectively serves as a preemptive warning system for potential maritime collisions.

The proposed solution was structured as a highly modular, six-stage software pipeline entirely driven by a Python backend. The detection tier relies upon the \gls{yolov8} neural network architecture, leveraging its anchor-free properties for rapid, scale-invariant vessel identification. This optical data is subsequently passed to the \gls{deepocsort} multi-object tracking framework. By integrating \gls{osnet} visual embeddings, the tracker learns the nuanced appearance features of individual ships, allowing the system to confidently re-identify targets even after they momentarily vanish behind coastal infrastructure or larger vessels. To bridge the gap between abstract pixel arrays and actionable geographic intelligence, the pipeline mathematically translates camera coordinates onto a standard geospatial plane. These coordinates are continuously fused with live \gls{ais} broadcast streams via a Haversine nearest-neighbor algorithm. The analytical tier then utilizes a \gls{bilstm} neural network to anticipate the likely path of each vessel, routing these forecasts directly into a Closest Point of Approach (\gls{cpa}) module designed to instantly alert surveillance operators to critical safety hazards.

Empirical evaluation against the demanding \gls{smd} dataset confirms that the MARVIS framework unequivocally satisfies all initial design parameters. The detection module returned a highly competitive mAP\textsubscript{50} of 0.91, validating its resilience against environmental fluctuations and severe variations in vessel size. During the tracking phase, the system achieved a \gls{mota} score of 87.4\% while recording an exceptionally low count of just 8 identity switches. In terms of sequence modeling, the \gls{lstm} trajectory predictor substantially outperformed standard kinematic baselines, realizing an Average Displacement Error (\gls{ade}) of 44.90 pixels—a 28.0\% error reduction that proved vital for mapping non-linear turning maneuvers. Supported by this predictive accuracy, the \gls{cpa} collision module attained a 0.92 precision rate when evaluating critical-risk encounters (under 50\,m). These cumulative metrics decisively prove that a sophisticated, deep-learning-powered maritime surveillance suite can be successfully constructed using open-source frameworks and effectively deployed on standard consumer-grade hardware.

% ─────────────────────────────────────────────────────────────────────────────
\section{Limitations}

Despite the strong quantitative performance achieved during testing, the current system architecture possesses specific operational limitations that must be acknowledged. Foremost, the accuracy of the geospatial projection is inextricably tied to the precision of the initial camera calibration. Any slight error in calculating the camera's tilt, pan, or lens distortion propagates mathematically, potentially skewing the projected geographic coordinates and thereby complicating the \gls{ais} sensor fusion process. Additionally, while the optical pipeline degrades gracefully, the system fundamentally relies on the visual spectrum. Consequently, detection reliability and overall operational efficacy can be noticeably compromised during extreme meteorological events, such as torrential rain or dense marine fog, which physically obscure the optical sensor.

\section{Future Enhancement}

% ─────────────────────────────────────────────────────────────────────────────

Although the MARVIS framework serves as a highly functional foundation, numerous avenues exist for expanding its capabilities and preparing the architecture for broader professional deployment.

A primary technical enhancement involves upgrading the foundational geospatial mathematics. The current methodology utilizes a linear field-of-view approximation, which performs reliably within a localized 2-kilometer surveillance radius. However, at extended ranges, this approach inevitably loses fidelity due to the geometric curvature of the Earth. Transitioning the logic to a full homographic transformation matrix anchored by precise \gls{gps} ground-control points would resolve this spatial degradation. Furthermore, if the system were adapted for aerial deployment via unmanned aerial vehicles, the projection algorithms could theoretically extract the necessary orientation variables directly from the drone's onboard telemetry, effectively eliminating the need for manual static calibration.

The collision forecasting logic also presents significant opportunities for advancement. At present, the \gls{cpa} calculations are governed by an assumption of constant linear velocity. A sophisticated upgrade would involve directly injecting the non-linear, 20-step curved path generated by the \gls{lstm} network straight into the collision detection matrix. Taking this algorithmic step would allow the system to anticipate disasters involving vessels that are currently maintaining a safe distance but are actively executing maneuvers that place them on a future collision trajectory.

From the perspective of artificial intelligence development, diversifying the training data remains a critical next step. The current trajectory prediction weights were optimized exclusively on traffic density patterns native to the Singapore Strait. Achieving global deployment readiness would necessitate rigorous transfer learning across diverse datasets gathered from vastly different maritime topologies. Moreover, anticipating human navigational behavior introduces inherent uncertainty. Evolving the predictive engine from a deterministic \gls{lstm} architecture to a probabilistic generative model—such as a Conditional Variational Autoencoder (CVAE) or a Generative Adversarial Network (GAN)—would enable the dashboard to output a probabilistic heat-map representing all possible future locations, rather than returning a single, absolute trajectory path.

Finally, pursuing hardware-level optimizations is essential for establishing true system autonomy. Transitioning the entire MARVIS suite onto a compact edge-computing platform, such as an \gls{nvidia} Jetson module, would require aggressive quantization of the deep learning weights down to \gls{int8} precision. Executing the visual inference processes directly at the edge, rather than transmitting video feeds to a centralized server, would drastically cut network latency and ensure that localized surveillance nodes remain fully operational even during severe communication blackouts.

% ─────────────────────────────────────────────────────────────────────────────
\section{Learning Outcomes}

% ─────────────────────────────────────────────────────────────────────────────

The comprehensive development of the MARVIS framework provided invaluable insights into the practical complexities of full-stack artificial intelligence engineering. The project necessitated moving beyond the theoretical isolation of a single machine learning algorithm to construct a robust, production-grade software architecture. This required designing a synchronized pipeline where deep computer vision inference, recurrent sequence modeling, and complex geospatial mathematics executed continuously and in parallel without causing computational bottlenecks.

Through the targeted fine-tuning of the \gls{yolov8} network, a deep structural understanding of anchor-free object detection paradigms was achieved. The practical process of balancing statistical confidence thresholds to suppress environmental false positives—such as those induced by aggressive sun glare or heavy marine chop—highlighted the nuanced differences between theoretical performance and real-world deployment. Additionally, integrating \gls{deepocsort} practically demonstrated the absolute necessity of appearance-based Re-Identification (\gls{reid}) embeddings for preserving target tracking continuity within densely populated environments.

The architecture and training of the \gls{bilstm} trajectory predictor offered advanced experience in applied sequence modeling. Engineering the custom loss function revealed that mathematical models require strict penalties for physical impossibilities, such as erratic and abrupt acceleration, to consistently generate realistic maritime paths. Furthermore, the challenge of mathematically aligning pixel-space observations with real-world \gls{ais} telemetry highlighted the substantial operational disparity that exists between optical camera feeds and standardized radar data.

Finally, networking the entire backend logic through asynchronous protocols, WebSockets, and a live interactive dashboard effectively bridged the technical gap between raw numerical data processing and the final user experience. Implementing these technologies fostered a deep understanding of how to manage asynchronous data streams while ensuring that heavy graphical rendering tasks did not stall the underlying visual inference pipeline. The systematic management of this extensive codebase through strict version control protocols ensured that the underlying logic remained perfectly stable throughout rapid evolutionary cycles. Ultimately, the successful end-to-end integration and rigorous evaluation of the proposed surveillance framework powerfully demonstrates its practical applicability and long-term effectiveness in the field of maritime security.